\begin{textsample}{NEG Dim 5 – human – Score -9.00 – mg/mg\_ef\_linguagens\_lp\_3\_p3.txt}  \label{ex:f5_neg_016}
Ao componente Língua Portuguesa cabe assegurar os direitos de aprendizagem aos \textbf{estudantes}, proporcionando -lhes experiências que contribuam para a ampliação e aprofundamento dos diferentes letramentos já adquiridos e aquisição de novos letramentos e multiletramentos, de forma a possibilitar a participação significativa e crítica nas diversas práticas sociais.
Os \textbf{objetos} de \textbf{conhecimento}, assim como as práticas de ensino, devem ser selecionados em função da aquisição e \textbf{desenvolvimento} das competências e \textbf{habilidades} de uso da língua e da reflexão sobre esse uso em diversas práticas, e não em função do domínio de conceitos e classificações como fins em si mesmos.
Assim, devem compor o currículo do componente \textbf{curricular} Língua Portuguesa aqueles conteúdos considerados \textbf{essenciais} à vida em sociedade, especialmente aqueles cuja aprendizagem exige intervenção e mediação sistemáticas da escola, para o \textbf{desenvolvimento} de competências específicas.
Vale destacar que essas competências perpassam todos os componentes \textbf{curriculares} do Ensino Fundamental e são \textbf{essenciais} para a ampliação das possibilidades de participação dos \textbf{estudantes} em práticas de diferentes campos de atividades humanas e de pleno exercício da cidadania.

% matched lemmas: conhecimento, curricular, desenvolvimento, essencial, estudante, habilidade, objeto
\end{textsample}
